\SourceLeaf{110}
\Article{18}
Suppose we place the second needle at different points on the same straight line, so that $\psi$ and $U$ stay fixed and only the distance $R$ changes, and observe the deflections of the movable needle from its equilibrium in the absence of the second needle, that is, the angles $u-u^0$. It is now clear that we can use elimination to find the values of as many of the coefficients $F,F',F''$ and so on as still have an appreciable effect. This gives us
\[
\frac MT=\left(1+\frac{\theta}{Tm}\right)\frac{F}{n\cos(\psi-U)\sin(\psi-u^0)+\sin(\psi-U)\cos(\psi-u^0)}
\]
Here the value of $\frac{\theta}{Tm}$ can be found by the method we described in Article $8$. The following observations will make the procedure easier in practice.

I. Instead of comparing $u$ with $u^0$, it is better to compare two opposite deflections by reversing the position of the second needle, that is, by keeping $R$ and $\psi$ fixed and increasing the angle $U$ by two right angles. Denote the corresponding values of $u$ by $u',u''$. With perfect symmetry we would have \textit{exactly} $u''=-u'$ if also $u^0=0$. But there is no need to adhere scrupulously to these conditions. Clearly, $u'$ and $u''$ are given by similar series whose first terms have \textit{exactly} opposite values. Thus $\tfrac12(u'-u'')$ and $\operatorname{tang}\tfrac12(u'-u'')$ are also given by a similar series whose first coefficient is \textit{exactly} $=F$.

II. It is better still to combine the experiments in groups of four, also changing the angle $\psi$ by two right angles, or taking the distance $R$ on the other side. If the values of $u$ in the latter two experiments are $u''',u''''$, the difference $\tfrac12(u'''-u'''')$ will likewise be expressed by a similar series, whose first term will also have coefficient $=F$. It is worth noting, as follows readily from the foregoing, that if $n$ were odd, the coefficients $F,F'',F''''$ and so on would be exactly equal to all orders in the two series for $u'-u^0$ and $u'''-u^0$, whereas the coefficients $F',F''',F^{\mathrm{v}}$ and so on would be exactly opposite to all orders. The same would hold for $u''-u^0$ and $u''''-u^0$, so that alternate terms would drop out of the series for $u'-u''+u'''-u''''$. But in nature, where $n=2$, this relation between the series for $u'-u^0$ and $u'''-u^0$ does not in general hold exactly. Already at the power $R^{-6}$, the coefficients fail to be exactly opposite. It can nevertheless be shown that this term too cancels completely in the combination $u'-u''+u'''-u''''$, so that $\operatorname{tang}\tfrac14(u'-u''+u'''-u'''')$ has the form
\[
LR^{-3}+L'R^{-5}+L''R^{-7}+\text{etc.}
\]
\SourceLeaf{111}or, more generally, if we leave the value of $n$ undetermined for now, the form
\[
LR^{-(n+1)}+L'R^{-(n+3)}+L''R^{-(n+5)}+\text{etc.}
\]
with $L=F$.

III. It is advisable to choose the angles $\psi,U$ so that small errors in measuring them do not appreciably change the value of $F$. To this end, the value of $U$ for a \textit{given} value of $\psi$ must be chosen so that $F$ attains a maximum. The condition is
\[
\operatorname{cotang}(\psi-U)=n\operatorname{tang}(\psi-u^0)
\]
which gives
\[
F=\pm\frac{mM}{mT+\theta}\BareRoot{(nn\sin(\psi-u^0)^2+\cos(\psi-u^0)^2)}
\]

The angle $\psi$, in turn, must be chosen so that this value of $F$ attains either a maximum or a minimum. The former occurs for $\psi-u^0=90^0$ or $270^0$, where
\[
F=\pm\frac{nmT}{mT+\theta}
\]
and the latter for $\psi-u^0=0$ or $180^0$, where
\[
F=\pm\frac{mM}{mT+\theta}
\]

